% GENERATED FILE. Edit canonical lessons, then run npm run generate:books.
% canonical-source-hash: fnv1a64:ec43340bea778bf6
% canonical-lessons: IT-C271-se, IT-C271-se-piove, IT-C271-dipende, IT-C271-ripasso-condizioni

\chapter{If It Rains --- Real Conditions, and It Depends}
\label{ch:it-real-conditions}
\hlchaptermodality{271}

\begin{chapteropening}
\textbf{By the end of this chapter:} \emph{I can set out a real condition with se and the present, and answer with dipende, naming what it depends on.}

Complete IT-C271-ripasso-condizioni: plans with a branch in them, from decided to not yet named.
\end{chapteropening}

\section[se]{se — if}

\label{lesson:IT-C271-se}



\begin{quote}
Say \textquotedblleft{}yes\textquotedblright{}. (\emph{Sì.}) Say \textquotedblleft{}I don't know\textquotedblright{}. (\emph{Non so.})

Every sentence so far has said what is, or what was. This chapter says what follows \textbf{if} something else is true, and it begins with two letters.
\end{quote}

\begin{grammarlens}[title={Grammar Lens: se opens a condition}]
\textbf{se} — \textquotedblleft{}if\textquotedblright{}.

\begin{quote}
\textbf{Se} sei stanco, dormi. — If you're tired, sleep. \textbf{Se} il treno è in ritardo, leggo. — If the train is late, I read.
\end{quote}

Whatever follows \emph{se} is the \textbf{condition}; the rest of the sentence is what follows from it. The condition can come first, with a comma, or second, without one: \emph{Leggo \textbf{se} il treno è in ritardo}.

\textbf{se} also means \textquotedblleft{}whether\textquotedblright{}, after a verb of knowing or asking, and here your \emph{non so} comes in useful:

\begin{quote}
\textbf{Non so se} è vero. — I don't know if it's true.
\end{quote}
\end{grammarlens}

\begin{cousinweb}
Listen to how close these are:

\begin{quote}
\textbf{se} — \textquotedblleft{}if\textquotedblright{} · \textbf{sì} — \textquotedblleft{}yes\textquotedblright{}, with its accent · \textbf{si} — \textquotedblleft{}oneself\textquotedblright{}, as in \emph{come si chiama?}
\end{quote}

\textbf{se} is inherited from Latin \emph{sī}, \textquotedblleft{}if\textquotedblright{} (EtymDB 2.1; Lewis \& Short, under \emph{si}); the vowel opened from \emph{i} to \emph{e}. \textbf{sì}, \textquotedblleft{}yes\textquotedblright{}, is a different Latin word: chapter twenty-nine traced it to \emph{sīc}, \textquotedblleft{}thus\textquotedblright{}. \emph{se} has its own vowel, so the ear can tell it from the other two. \emph{sì} and \emph{si} sound exactly alike, and the accent on \emph{sì} is there to keep those two apart on the page, where the ear cannot help.
\end{cousinweb}

\subsection*{Your turn}
Say these aloud:

\begin{itemize}
\raggedright
  \item \emph{Se sei stanco, dormi.}
  \item the same, condition second — \emph{Dormi se sei stanco.}
  \item \emph{Se il treno è in ritardo, leggo.}
  \item \emph{Non so se è vero.}
\end{itemize}

\begin{tcolorbox}[breakable,colback=teal!4,colframe=teal!35!black,title={Before you move on}]
Which half of \emph{se sei stanco, dormi} is the condition? (\emph{\textbf{Se sei stanco.}}) What does \emph{se} mean after \emph{non so}? (\textbf{Whether}, \textquotedblleft{}if\textquotedblright{}.) Which Latin words are behind \emph{se} and \emph{sì}? (\emph{\textbf{Sī}}, \textquotedblleft{}if\textquotedblright{}, and \emph{\textbf{sīc}}, \textquotedblleft{}thus\textquotedblright{}.)
\end{tcolorbox}
\section[se piove, resto a casa]{se piove, resto a casa — the condition that really can happen}

\label{lesson:IT-C271-se-piove}



\begin{quote}
Say \textquotedblleft{}if\textquotedblright{}. (\emph{Se.}) Say \textquotedblleft{}it was raining\textquotedblright{}. (\emph{Pioveva.}) Now the present: \textquotedblleft{}it's raining\textquotedblright{}. (\emph{Piove.})
\end{quote}

\begin{grammarlens}[title={Grammar Lens: both halves in the present}]
\begin{quote}
\textbf{Se piove, resto a casa.} — If it rains, I stay at home. \textbf{Se il treno è in ritardo, prendo l'autobus.} — If the train is late, I take the bus.
\end{quote}

This is a \textbf{real} condition: it may well happen, and if it does, the result follows. Italian puts both halves in the plain present, as English does. There is nothing new to build. \emph{se} plus a sentence you could already say, then another sentence you could already say.

The result half can also be advice, with the \emph{dovresti} of chapter two hundred and sixty-six:

\begin{quote}
\emph{Se sei stanco, \textbf{dovresti} dormire.} — If you're tired, you should sleep.
\end{quote}

And it can be a plain instruction, as in the last lesson: \emph{se sei stanco, dormi}.
\end{grammarlens}

\begin{culture}
A real condition is a plan with a branch in it. Travel is full of them, and the station gives you two at once:

\begin{quote}
\emph{Se perdo il treno, prendo l'autobus. Se perdo l'autobus, resto qui.}
\end{quote}

Compare \emph{quindi}. \emph{Piove, quindi resto a casa} says it \textbf{is} raining, and here is the consequence. \emph{Se piove, resto a casa} says it \textbf{might}, and here is what happens if so. Same consequence; the second one has not decided yet whether the cause is true.
\end{culture}

\begin{cousinweb}
\textbf{piovere} is inherited from spoken Latin \emph{plŏvere} (EtymDB 2.1), a reshaping of classical \emph{pluere}, \textquotedblleft{}to rain\textquotedblright{} (Lewis \& Short, under \emph{pluo}). The Latin \emph{pl} became Italian \emph{pi}, as in \textbf{più} from \emph{plūs} and \textbf{piegare} from \emph{plicāre}: three words, one sound change.
\end{cousinweb}

\subsection*{Your turn}
Say these aloud:

\begin{itemize}
\raggedright
  \item \emph{Se piove, resto a casa.}
  \item \emph{Se il treno è in ritardo, prendo l'autobus.}
  \item the branch — \emph{Se perdo il treno, prendo l'autobus.}
  \item advice in the result — \emph{Se sei stanco, dovresti dormire.}
\end{itemize}

\begin{tcolorbox}[breakable,colback=teal!4,colframe=teal!35!black,title={Before you move on}]
Which tense goes in each half of a real condition? (The \textbf{present}, in both.) What is the difference between \emph{piove, quindi resto a casa} and \emph{se piove, resto a casa}? (The first says it \textbf{is} raining; the second has \textbf{not decided}.) Which other words show Latin \emph{pl} becoming \emph{pi}? (\emph{\textbf{Più}} and \emph{\textbf{piegare}}.)
\end{tcolorbox}
\section[dipende]{dipende — the condition you do not have to finish}

\label{lesson:IT-C271-dipende}



\begin{quote}
Say \textquotedblleft{}if it rains, I stay at home\textquotedblright{}. (\emph{Se piove, resto a casa.})

That sentence names its condition. Sometimes you want the condition without saying what it is.
\end{quote}

\subsection*{You'll want to know: dipende}
\textbf{dipende} — \textquotedblleft{}it depends\textquotedblright{}.

\begin{quote}
— \emph{Vai alla festa domani?} — \emph{Dipende.} — Are you going to the party tomorrow? — It depends.
\end{quote}

You have \textbf{dipendere}, \textquotedblleft{}to depend\textquotedblright{}. \emph{Dipende} is its \textquotedblleft{}it\textquotedblright{} form, and on its own it is a complete answer: it says a condition exists and declines to say which. That is exactly why it is so useful to a question you are not ready to answer.

\begin{grammarlens}[title={Grammar Lens: dipende da}]
Add \textbf{da} and you can name the condition after all:

\begin{quote}
\textbf{Dipende da} te. — It depends on you. \textbf{Dipende dal} treno. — It depends on the train.
\end{quote}

\emph{da} + \emph{il} is \textbf{dal}, the same welding as \emph{al} and \emph{del}. English says \emph{depend \textbf{on}}; Italian says \emph{dipende \textbf{da}}, \textquotedblleft{}hangs \textbf{from}\textquotedblright{}, which is the older picture.

Or name it with a whole \emph{se}: \emph{Dipende. Se piove, resto a casa.} The \emph{dipende} announces a condition; the \emph{se} sentence states it.
\end{grammarlens}

\begin{cousinweb}
\textbf{dipendere} comes from Latin \emph{dēpendēre}, \textquotedblleft{}to hang from, to hang down\textquotedblright{} (Lewis \& Short, under \emph{dependeo}): \emph{dē}, \textquotedblleft{}down from\textquotedblright{}, and \emph{pendēre}, \textquotedblleft{}to hang\textquotedblright{}. To depend on something is to \textbf{hang from} it, and the answer swings until the condition settles. English \textbf{pendant} (a thing that hangs), \textbf{pending} (still hanging) and \textbf{suspend} share the root.
\end{cousinweb}

\subsection*{Your turn}
Say these aloud:

\begin{itemize}
\raggedright
  \item answer \emph{Vai alla festa domani?} — \emph{Dipende.}
  \item name it — \emph{Dipende dal treno.}
  \item name it with a whole condition — \emph{Dipende. Se piove, resto a casa.}
  \item hand it back — \emph{Dipende da te. Secondo te?}
\end{itemize}

\begin{tcolorbox}[breakable,colback=teal!4,colframe=teal!35!black,title={Before you move on}]
What does \emph{dipende} say, and what does it not say? (That there \textbf{is} a condition — not \textbf{which}.) What do you add to name it? (\emph{\textbf{Da}}: \emph{dipende da te}.) What is hanging, in \emph{dēpendēre}? (\textbf{The answer}, from the condition.)
\end{tcolorbox}
\section[(review)]{Review — plans with a branch in them}

\label{lesson:IT-C271-ripasso-condizioni}



\begin{quote}
Someone asks \emph{Vai alla festa domani?} Give the answer that decides nothing. (\emph{Dipende.})
\end{quote}

\begin{grammarlens}[title={What you've built}]
Three ways now to talk about a cause, and they differ in how sure you are.

\begin{quote}
\emph{Piove, \textbf{quindi} resto a casa.} — It is raining: decided. \emph{\textbf{Se} piove, resto a casa.} — It might: the branch is ready. \emph{\textbf{Dipende.}} — Not even the branch is named yet.
\end{quote}

Put them together at a ticket window and you have a whole conversation:

\begin{quote}
— \emph{Vorrei un biglietto per Roma.} — \emph{Andata e ritorno?} — \emph{Dipende. Se torno domani, andata e ritorno. Se resto, solo andata.}
\end{quote}
\end{grammarlens}

\subsection*{Your turn}
Say these aloud:

\begin{itemize}
\raggedright
  \item decided — \emph{Piove, quindi resto a casa.}
  \item a branch — \emph{Se piove, resto a casa.}
  \item not even the branch — \emph{Dipende.} — then name what it hangs from — \emph{Dipende dal treno.}
  \item at the window — \emph{Dipende. Se torno domani, andata e ritorno. Se resto, solo andata.}
  \item a plan of your own, with a \emph{se}: what you do if it rains this weekend
\end{itemize}

\begin{tcolorbox}[breakable,colback=teal!4,colframe=teal!35!black,title={Before you move on}]
Order the three from most sure to least. (\emph{\textbf{Quindi}}, then \emph{\textbf{se}}, then \emph{\textbf{dipende}}.) Which tense goes in both halves of \emph{se piove, resto a casa}? (The \textbf{present}.) Which Latin word is \emph{se}? (\emph{\textbf{Sī}}, \textquotedblleft{}if\textquotedblright{}.)
\end{tcolorbox}
